
\documentclass[10pt,a4paper]{article}
 
\usepackage[T1]{fontenc}
\usepackage[utf8]{inputenc}
\usepackage[scaled=0.92]{helvet}
\renewcommand{\familydefault}{\sfdefault}


\usepackage[a4paper,left=1.6cm,right=1.6cm,top=0.9cm,bottom=0.8cm]{geometry}
\usepackage{xcolor}
\usepackage{graphicx}
\usepackage{tikz}
\usepackage{paracol}
\usepackage[most]{tcolorbox}
\usepackage{enumitem}
\usepackage[hidelinks]{hyperref}
 
\pagestyle{empty}
\setlength{\parindent}{0pt}
\setlength{\columnsep}{1.0cm}
 
% ---------------------------------------------------------------------
%  Farben
% ---------------------------------------------------------------------
\definecolor{accent}{HTML}{2E7DD1}
\definecolor{linkblue}{HTML}{1A0DAB}
\definecolor{softgray}{HTML}{8A8A8A}
\definecolor{tagbg}{HTML}{ECECEC}
\definecolor{barbg}{HTML}{E4F1FA}
\definecolor{dotoff}{HTML}{C8C8C8}
 
% ---------------------------------------------------------------------
%  Icons
% ---------------------------------------------------------------------
\usepackage{marvosym}
\newcommand{\faPhone}{\Mobilefone}
\newcommand{\faEnvelope}{\Letter}
\newcommand{\faGlobe}{\ComputerMouse}
\newcommand{\faLinkedin}{\textbf{in}}
\newcommand{\faGithub}{\textbf{GH}}
\newcommand{\faMapMarker}{\Mundus}
 
% ---------------------------------------------------------------------
%  Bausteine
% ---------------------------------------------------------------------
% Abschnittsüberschrift mit Linie
\newcommand{\cvsection}[1]{%
  \vspace{7pt}%
  {\large\bfseries\MakeUppercase{#1}}\\[-7pt]
  \rule{\linewidth}{0.6pt}\par\vspace{4pt}}
 
% Eintrag: {Titel}{Institution}{Zeitraum}{Ort}
\newcommand{\cventry}[4]{%
  {\bfseries #1}\par
  {\itshape\color{black!75}#2}\par
  {\small\color{softgray}#3\hspace{1.2em}#4}\par\vspace{3pt}}
 
% Aufzählung
\newenvironment{cvitems}{%
  \begin{itemize}[leftmargin=1.2em,topsep=3pt,itemsep=2pt,parsep=0pt,label=\textbullet]}%
  {\end{itemize}\vspace{4pt}}
 
% Blauer Balken für Skill-Kategorien
\newcommand{\skillhead}[1]{%
  \par\vspace{2pt}%
  \colorbox{barbg}{\parbox{\dimexpr\linewidth-2\fboxsep}{\bfseries #1}}\par\vspace{3pt}}
 
% Graue "Pill"-Tags
\newtcbox{\skilltag}{on line, colback=tagbg, colframe=tagbg, arc=6pt,
  boxsep=0pt, left=6pt, right=6pt, top=2.5pt, bottom=2.5pt,
  boxrule=0pt, fontupper=\small\bfseries\color{black!80}}
 
% Sprach-Punkte: \langdots{gefüllt}{gesamt}
\newcommand{\langdots}[2]{%
  \foreach \i in {1,...,#2}{%
    \ifnum\i>#1 {\color{dotoff}\textbullet}\else{\color{accent}\textbullet}\fi}}
 
% Link
\newcommand{\cvlink}[2]{\href{#1}{\color{linkblue}#2}}
 
% Foto: foto.jpeg im Projektordner, wird rund zugeschnitten
% Ohne Datei erscheint ein grauer Platzhalter.
\newcommand{\cvphoto}{%
  \begin{tikzpicture}
    \IfFileExists{foto.jpeg}{%
      \begin{scope}
        \clip (0,0) circle (1.25cm);
        \node at (0,-0.30) {\includegraphics[width=4.2cm]{foto.jpeg}};
      \end{scope}
      \draw[accent, line width=1.2pt] (0,0) circle (1.25cm);
    }{%
      \fill[accent] (0,0) circle (1.25cm);
      \node[text=white,font=\fontsize{26}{28}\selectfont\bfseries] at (0,0) {BZ};
    }
  \end{tikzpicture}}
 
% =====================================================================
\begin{document}
 
% --------------------------- KOPFZEILE -------------------------------
\begin{minipage}[c]{3.0cm}
  \cvphoto
\end{minipage}%
\hspace{0.6cm}%
\begin{minipage}[c]{\dimexpr\linewidth-3.6cm}
  {\fontsize{28}{32}\selectfont\MakeUppercase{Abo Bakir Fawaz Zanoun}}\par\vspace{4pt}
  {\Large\color{accent}Computer Engineer (M.Eng.) \textbar{} Software, DevOps \& IoT}\par\vspace{6pt}
  \small
  \faPhone~0157\,/\,32236600\hspace{1.2em}
  \faEnvelope~\cvlink{mailto:bakirzanoun@gmail.com}{bakirzanoun@gmail.com}\hspace{1.2em}
  \faGlobe~\cvlink{https://bakirzanoun.de}{bakirzanoun.de}
  \par\vspace{3pt}
    \faLinkedin~\cvlink{https://www.linkedin.com/in/bakir-zanoun-226b09108}{linkedin.com/in/bakir-zanoun}
  \hspace{1.2em}
  \faGithub~\cvlink{https://github.com/Bakir-Alnuaimi}{github.com/Bakir-Alnuaimi}
  \par\vspace{3pt}
  {\color{black!75}\faMapMarker~Berlin, Deutschland\hspace{1.2em}Staatsangehörigkeit: Deutsch}
\end{minipage}
 
\vspace{6pt}
 
% --------------------------- ZWEI SPALTEN ----------------------------
\columnratio{0.5}
\begin{paracol}{2}
 
% ============================ LINKS ==================================
\cvsection{Ausbildung}
 
\cventry{Master of Engineering (M.Eng.) in Computer Engineering}
        {HTW Berlin – Hochschule für Technik und Wirtschaft}
        {10/2024 – 09/2026}{Berlin}
\begin{cvitems}
  \item Abschluss am 22.09.2026 \textbar{} Masterarbeit \textbf{auf Englisch}
\end{cvitems}
 
\cventry{Bachelor of Engineering (B.Eng.) in Computer Engineering}
        {HTW Berlin – Hochschule für Technik und Wirtschaft}
        {10/2021 – 03/2024}{Berlin}
\begin{cvitems}
  \item Bachelorarbeit \textbf{auf Deutsch} (F\&E-Projekt bei AUCOTEAM, Trinkwasserüberwachung)
\end{cvitems}
 
\cventry{Informationstechnischer Assistent + Fachabitur}
        {Staatlich geprüft -- AUCOTEAM Berufsfachschule, Berlin-Pankow}
        {09/2018 – 08/2021}{Berlin}
 
\cvsection{Berufserfahrung}
 
\cventry{Werkstudent Softwareentwicklung}
        {AUCOTEAM GmbH}
        {07/2023 – 09/2026}{Berlin}
\begin{cvitems}
  \item GitLab-CI/CD-Pipeline (DevSecOps) für Container kon\-zi\-piert und umgesetzt: Test, Build, Trivy-Scan, Push, Teams-Alarm
  \item Zeitgesteuerte Re-Scans der Registry auf neue CVEs für zwei produktive .NET-Services
  \item Multisensor-Prüfstand im Förderprojekt zur Trink\-wasser\-über\-wachung (Raspberry Pi, I2C, Python, InfluxDB)
  \item F\&E-Projekte zu Embedded Systems und IoT-Lösungen mit LoRaWAN und MQTT
  \item Softwareentwicklung mit .NET, C\# und ASP.NET; KI-Chatbot inkl. API zwischen Frontend und Backend
\end{cvitems}
 
\cventry{Helfer in der Flüchtlingshilfe}
        {Johanniter-Unfall-Hilfe – Ukraine-Ankunftszentrum Tegel}
        {03/2022 – 12/2022}{Berlin}
 
\cventry{Call-Center-Agent (Identitätsprüfung)}
        {hello.de für IDnow}
        {09/2019 – 02/2020}{Berlin}
\begin{cvitems}
  \item Identitätsprüfung per Video-Ident: Echtheitsprüfung von Ausweisdokumenten aus aller Welt
  \item Erkennung von Betrugsversuchen und Social-Engineering-Angriffen
\end{cvitems}
 
\cventry{Aushilfskraft Lagerlogistik}
        {Zalando GmbH}
        {10/2017 – 02/2018}{Erfurt}
{\small Konfektionierung, IT-gestützte Lagerverwaltung}\par\vspace{4pt}
 
\cventry{Krankenpflegepraktikum}
        {HELIOS Klinikum Gotha}
        {04/2017}{Gotha}
{\small Grundpflege, Pflegedokumentation, Übersetzung (DE–AR)}\par
 
% ============================ RECHTS =================================
\switchcolumn
 
\cvsection{Technische Kenntnisse}
{\raggedright
 
\skillhead{Programmiersprachen:}
\skilltag{C\#} \skilltag{Python} \skilltag{C} \skilltag{C++} \skilltag{Bash}
 
\skillhead{Frameworks \& Web:}
\skilltag{.NET} \skilltag{ASP.NET} \skilltag{REST API} \skilltag{HTML/CSS/JS}
 
\skillhead{DevOps \& Security:}
\skilltag{GitLab CI/CD} \skilltag{GitLab Runner} \skilltag{Trivy} \skilltag{Docker} \skilltag{Docker Compose} \skilltag{OCI Registry}
 
\skillhead{Embedded \& IoT:}
\skilltag{Raspberry Pi} \skilltag{ESP32} \skilltag{STM32} \skilltag{Arduino} \skilltag{Jetson Nano} \skilltag{ROS} \skilltag{I2C} \skilltag{SPI} \skilltag{UART} \skilltag{MQTT} \skilltag{LoRaWAN} \skilltag{InfluxDB}
 
\skillhead{Systeme \& Tools:}
\skilltag{Linux/Ubuntu} \skilltag{Netzwerke} \skilltag{Jira} \skilltag{Git} \skilltag{KI / Chatbots} \skilltag{AutoCAD}\par}
 
\cvsection{Abschlussarbeiten}
 
{\bfseries Masterarbeit} {\small\color{softgray}(verfasst auf Englisch)}\par
\textit{Designing and Implementing a CI/CD Pipeline Architecture with Continuous Automated Security Scans for Containerized Applications}\par
\vspace{8pt}
 
{\bfseries Bachelorarbeit} {\small\color{softgray}(verfasst auf Deutsch)}\par
\textit{Entwicklung und Implementierung einer Multisensorphalanx für schnelle Integrationen diverser Sensoren in einem Testaufbau}\par
 
\cvsection{Zertifikate}
 
\begin{cvitems}
  \item \textbf{Einführung EU-DSGVO} – IITR Datenschutz GmbH {\color{softgray}(10/2019)}
  \item \textbf{Basic Privacy Training} – IITR Datenschutz GmbH {\color{softgray}(10/2019)}
  \item \textbf{Führerschein Klasse B}
\end{cvitems}
 
\cvsection{Sprachen}
\renewcommand{\arraystretch}{1.15}
\begin{tabular}{@{}l@{\hspace{1em}}l@{\hspace{0.6em}}l@{}}
  \textbf{Arabisch:}    & \langdots{5}{5} & Muttersprache \\
  \textbf{Deutsch:}     & \langdots{5}{5} & Fließend \\
  \textbf{Englisch:}    & \langdots{4}{5} & Fließend \\
  \textbf{Französisch:} & \langdots{1}{5} & Grundkenntnisse \\
\end{tabular}
 
\cvsection{Ehrenamt \& Interessen}
\textbf{Ehrenamtlicher Übersetzer} (Deutsch – Arabisch)\par\vspace{3pt}
Schwimmen, Musizieren und Gitarre spielen
 
\vspace{12pt}
{\small\color{softgray}Berlin, September 2026}
 
\end{paracol}
 
\end{document}
